\section{The Topology of the Prime Pairing}
\label{v4-arith:w6-b:union-density:closure}

The critical prime measure has exponential growth in the logarithmic
variable. Polynomial decay of a test function does not control its
prime sum. The convolution square already detects this failure.

\begin{lemma}[the proposed polynomial majorant diverges]
For every real $N$,
\[
 \sum_p\sum_{m\ge1}(\log p)p^{-m/2}(m\log p)^{-N}=+\infty.
\]
\end{lemma}
\begin{proof}
Its $m=1$ summand is $p^{-1/2}(\log p)^{1-N}\ge p^{-1}$ for
sufficiently large primes. The reciprocal prime sum diverges.
Indeed convergence of $\sum_p1/p$ would bound
$\prod_p(1-1/p)^{-1}$, since $-\log(1-1/p)\le2/p$ for $p\ge2$.
The finite Euler products contain every integer reciprocal with
prime factors in that finite set. Their limit therefore dominates
all harmonic partial sums, which are unbounded.
\end{proof}

\begin{theorem}[failure on a positive Schwartz square]
\label{v4-arith:top:schwartz-counterexample}
Let $f(t)=\exp(-(1+t^2)^{1/4})$.
Then $f\in\mathscr S(\mathbb R)$, but
\[
 \sum_{n\ge2}\frac{\Lambda(n)}{\sqrt n}
 \bigl((f*f^*)(\log n)+(f*f^*)(-\log n)\bigr)=+\infty.
\]
There are nonnegative $f_j\in\mathscr A$ with $f_j\to f$ in
$\mathscr S$ and $W_{\mathrm{fin}}(f_j*f_j^*)\to+\infty$.
Consequently the compact prime functional has no continuous
extension to $\mathscr S$ with these values.
\end{theorem}
\begin{proof}
Every derivative of $f$ is a finite sum of products of derivatives
of $(1+t^2)^{1/4}$ times $f$. These products have polynomial growth,
whereas $f$ decays faster than every inverse power. Thus $f$ is Schwartz.
For $t\ge2$, integrate the autocorrelation over $v\in[0,1]$ to obtain
\[
 (f*f^*)(t)=\int f(v)f(v-t)dv\ge c e^{-C\sqrt t}
\]
for constants $c,C>0$.
At $t=\log p$, the weighted term is eventually at least $p^{-1}$,
because $c(\log p)e^{(\log p)/2-C\sqrt{\log p}}\to\infty$.
The preceding lemma proves divergence using only the prime subseries.

Choose an even smooth cutoff $\chi$ with $0\le\chi\le1$, equal to
one on $[-1,1]$, zero off $[-2,2]$, and decreasing with $|t|$.
Put $f_j(t)=\chi(t/j)f(t)$.
Leibniz's rule and the Schwartz tails give $f_j\to f$ in every
Schwartz seminorm. The sequence increases pointwise to $f$.
Its nonnegative autocorrelations increase to $f*f^*$ by monotone
convergence. A second use of monotone convergence on the prime sum
gives the asserted divergence. A continuous extension would instead
have a finite limit, since convolution is continuous on $\mathscr S$.
\end{proof}

\begin{theorem}[the actual dominated and exponential limits]
\label{v4-arith:w6-b:thm:schwartz-closure}
Let $\varphi_j,\varphi\in\mathscr S$ with
$\varphi_j\to\varphi$ in $\mathscr S$.
If a nonnegative sequence $(d_n)$ satisfies
\[
 |\varphi_j(\pm\log n)|\le d_n\quad\hbox{for every }j,n,
 \qquad \sum_{n\ge2}\frac{\Lambda(n)}{\sqrt n}d_n<\infty,
\]
then the absolutely convergent prime pairings satisfy
$W_{\mathrm{fin}}(\varphi_j)\to W_{\mathrm{fin}}(\varphi)$.
In particular this convergence holds when $f_j\to f$ in the
rapid exponential space $\mathscr G$ and
$\varphi_j=f_j*f_j^*$, $\varphi=f*f^*$.
It then holds for the complete Weil form as well.
\end{theorem}
\begin{proof}
Point evaluation is continuous in the Schwartz topology.
The assumed summable sequence permits dominated convergence in the
prime sum, including both signs of the logarithm.
For the exponential statement, convolution continuity in
Proposition~\ref{v4-arith:sc:gaussian-extension} gives
$\varphi_j\to\varphi$ in $\mathscr G$.
A convergent sequence is bounded in the seminorm
$p_{2,0,0}(\varphi)=\sup_t e^{2|t|}|\varphi(t)|$.
Thus one may take $d_n=Cn^{-2}$, for which the sum is bounded by
$C\sum_{n\ge2}(\log n)n^{-5/2}<\infty$.
The same proposition controls the polar, gamma, and zero terms,
so the full explicit formula and its quadratic pairing pass to the limit.
\end{proof}

A norm identity on a dense subspace extends only in a topology that
controls both sides. A positive automorphic norm cannot supply the
missing continuity of a different prime distribution.
The actual signed finite-prime and full-zero comparisons remain
\eqref{v4-arith:sc:prime-comparison} and
\eqref{v4-arith:sc:full-comparison} on their stated carriers.
